\chapter{Ferramentas de projeto}
\label{cap:fundamentos}

Este capítulo reúne as ferramentas usadas em todos os controladores: especificações de
resposta, alocação de polos, lugar geométrico das raízes, margens de estabilidade e
controle ótimo.

\section{Especificações de um sistema de 2ª ordem}
Quase todas as leis lineares deste trabalho levam o erro de um eixo a uma dinâmica de
segunda ordem
\begin{equation}
  s^2 + 2\zeta\omega_n s + \omega_n^2 = 0 ,
  \qquad s_{1,2} = -\zeta\omega_n \pm j\omega_n\sqrt{1-\zeta^2}.
\end{equation}
\begin{table}[H]
\centering
\begin{tabular}{@{}lll@{}}
\toprule
Especificação & Fórmula & Inversa (projeto) \\
\midrule
Sobressinal & $M_p = e^{-\pi\zeta/\sqrt{1-\zeta^2}}$ &
  $\zeta = \dfrac{-\ln M_p}{\sqrt{\pi^2+\ln^2 M_p}}$ \\
Tempo de acomodação (2\%) & $t_s \approx 4/(\zeta\omega_n)$ & $\omega_n = 4/(\zeta t_s)$ \\
Tempo de subida (10--90\%) & $t_r \approx (1{,}8)/\omega_n$ & $\omega_n \approx 1{,}8/t_r$ \\
Tempo de pico & $t_p = \pi/(\omega_n\sqrt{1-\zeta^2})$ & --- \\
\bottomrule
\end{tabular}
\end{table}

\begin{receita}
\textbf{Especificar a perseguição.} Exemplo: ``recuperar um erro de distância de
\SI{1}{m} em \SI{2,5}{s} sem sobressinal visível''. Tome $M_p\le 5\%\Rightarrow\zeta\approx0{,}7$
e $t_s=\SI{2,5}{s}\Rightarrow\omega_n = 4/(0{,}7\cdot 2{,}5)=\SI{2,3}{rad/s}$.
Verifique depois se $\omega_n$ cabe na banda permitida pela PX4 e pelo atraso
(\cref{sec:regra-banda}).
\end{receita}

\section{Alocação de polos no duplo integrador}
Para $G(s)=1/s^2$ e a lei $u=-k_p e_p - k_d e_v$:
\begin{equation}
  s^2 + k_d s + k_p = 0
  \quad\Longrightarrow\quad
  \boxed{k_p = \omega_n^2,\qquad k_d = 2\zeta\omega_n.}
  \label{eq:pd-polos}
\end{equation}
Com $\omega_n=2{,}3$ e $\zeta=0{,}7$: $k_p=5{,}3$, $k_d=3{,}2$. Em MATLAB:
\code{K = place(A,B,p)} (contínuo) ou \code{K = place(Ad,Bd,exp(p*Ts))} (discreto).

\section{Lugar geométrico das raízes (LGR)}
Para uma malha $1 + k\,L(s) = 0$, o LGR é o conjunto das raízes quando $k$ varia de 0
a $\infty$. As regras básicas:
\begin{enumerate}
  \item Começa nos polos de $L$ ($k=0$) e termina nos zeros ($k\to\infty$) ou no infinito.
  \item No eixo real, pertencem ao LGR os trechos à esquerda de um número ímpar de
        polos+zeros reais.
  \item Assíntotas: $n-m$ retas com ângulos $(2q+1)180^\circ/(n-m)$ centradas em
        $\sigma_a = (\sum p_i - \sum z_i)/(n-m)$.
  \item Condição de ângulo: $\angle L(s) = 180^\circ(2q+1)$. Condição de módulo:
        $k = 1/|L(s)|$. A segunda dá o ganho que leva um polo a um ponto desejado.
\end{enumerate}

\paragraph{Como usar para calcular ganho.} Marque no plano $s$ (ou $z$) a região
permitida pelas especificações: $\zeta\ge\zeta_{min}$ (setor) e
$\zeta\omega_n\ge 4/t_s$ (à esquerda de uma vertical). O ganho é o valor de $k$ em que o
ramo dominante entra nessa região. No MATLAB:
\begin{lstlisting}
L = tf(1,[1 0 0]) * tf([1 a],1);   % ex.: PD  kd(s + kp/kd)
rlocus(L); sgrid(0.7, 2.3);        % linhas de zeta e wn
[k, polos] = rlocfind(L);          % clique no ponto desejado
\end{lstlisting}
No discreto, use \code{rlocus(Lz)} com \code{zgrid}. No plano $z$, a região estável é o
círculo unitário (\cref{cap:discreto}).

\section{Resposta em frequência e margens}
Com a malha aberta $L(j\omega)$, a frequência de cruzamento $\omega_c$ é aquela em que
$|L|=1$. A \textbf{margem de fase} é $\text{PM} = 180^\circ + \angle L(j\omega_c)$.
Valores usuais: PM $\ge 45^\circ$ (sobressinal $\lesssim 25\%$). Como regra,
$\zeta\approx \text{PM}/100$.

\paragraph{Atraso.} Um atraso puro $e^{-s\tau}$ não muda o módulo, mas subtrai
$\omega\tau$ radianos de fase. A \textbf{margem de atraso} é o atraso extra que
instabiliza a malha:
\begin{equation}
  \tau_{max} = \frac{\text{PM [rad]}}{\omega_c}.
  \label{eq:margem-atraso}
\end{equation}
\label{sec:regra-banda}
Daí sai a regra prática mais importante deste trabalho:
\begin{equation}
  \boxed{\omega_c \lesssim \frac{\text{PM}_{\text{livre}} - \text{PM}_{\text{desejada}}}{\tau_{total}}}
\end{equation}
Com $\tau_{total}\approx\SI{0,1}{s}$ e \ang{20}--\ang{30} de margem livre para gastar,
$\omega_c\lesssim\SI{3,5}{rad/s}$. Acima disso, o atraso da câmera e do rastreador
domina e a malha oscila.

No MATLAB: \code{margin(L)}, \code{[Gm,Pm,Wcg,Wcp] = margin(L)},
\code{L.InputDelay = tau} e \code{allmargin(L)} (inclui \code{DelayMargin}).

\section{Índices de desempenho}
Para comparar e otimizar controladores, use um índice escalar:
\begin{align}
  \text{ISE} &= \int_0^{T} e^2\,dt, &
  \text{IAE} &= \int_0^{T} |e|\,dt, &
  \text{ITAE} &= \int_0^{T} t\,|e|\,dt, \\
  J_{LQ} &= \int_0^{T} \left(x\T Q x + u\T R u\right)dt . &&&&
  \label{eq:jlq}
\end{align}
O ITAE penaliza erros que persistem e dá respostas bem amortecidas. O $J_{LQ}$ é o
índice do LQR/MPC e equilibra erro e esforço.

\section{Controle ótimo: visão geral}
\label{sec:otimo}
``Controle ótimo'' quer dizer escolher a lei que minimiza um custo. Neste trabalho há
três formas, usadas em capítulos diferentes:
\begin{enumerate}
  \item \textbf{Ótimo analítico (LQR).} Para modelo linear e custo quadrático
        \eqref{eq:jlq} com horizonte infinito, a solução é $u=-Kx$, com $K$ vindo da
        equação de Riccati. Tempo contínuo:
        \begin{equation}
          A\T P + PA - PBR^{-1}B\T P + Q = 0,\qquad K = R^{-1}B\T P .
        \end{equation}
        Tempo discreto (DARE, usada no código):
        \begin{equation}
          P = Q + A\T P A - A\T P B (R + B\T P B)^{-1} B\T P A,\qquad
          K = (R + B\T P B)^{-1} B\T P A .
          \label{eq:dare}
        \end{equation}
  \item \textbf{Ótimo com restrições (MPC).} Minimiza o mesmo custo num horizonte finito,
        com limites de $|u|$ e $|v|$, resolvendo um problema quadrático (QP) a cada ciclo.
  \item \textbf{Ajuste ótimo de parâmetros.} Para leis que não são ótimas por
        construção (PID, SMC, guiamento), fixa-se a estrutura e busca-se os ganhos
        $\theta$ que minimizam $J(\theta)$ \emph{simulando} a malha fechada completa
        (com PX4, atraso, ruído):
        \begin{equation}
          \theta^\star = \arg\min_\theta\; J(\theta)
          \quad\text{(\code{fminsearch}, \code{patternsearch}, \code{ga}).}
        \end{equation}
        É a forma mais justa de comparar técnicas diferentes: cada uma com o melhor
        ajuste possível para o \emph{mesmo} custo e os \emph{mesmos} cenários
        (\cref{cap:comparacao}).
\end{enumerate}
